% Simulazioni d'esame. I quesiti vengono da esercizi.toml (ADR 0006): si modificano lì.
\simulazioni

Ogni simulazione è una prova completa, come quella vera: 31 quesiti su tutte le unità didattiche, 21 a risposta multipla e 10 a completamento, da svolgere in 50 minuti. I quesiti sono divisi fra le unità didattiche in proporzione al loro peso nel syllabus.

\begin{tcolorbox}[fisbox={formula}{formulaBg},title={Come usare le simulazioni}]
Fai una simulazione solo quando hai studiato tutte le unità didattiche. Mettiti al tavolo con una penna, due fogli bianchi per i conti e un orologio: niente calcolatrice, niente libro, niente appunti. Hai 50 minuti.

Nei quesiti a risposta multipla una sola delle cinque opzioni è corretta. Nei quesiti a completamento scrivi un numero o una parola, al massimo 16 caratteri, con la virgola come separatore decimale e senza separatore delle migliaia; l'unità di misura è già nel testo.

Alla fine correggi con la tabella delle risposte: \textbf{+1} per ogni risposta esatta, \textbf{0} per ogni risposta non data, \textbf{$-0{,}1$} per ogni risposta sbagliata. La prova è superata con almeno 18 punti su 31. Poi leggi le soluzioni svolte dei quesiti sbagliati: accanto a ogni quesito c'è l'unità didattica da ripassare.
\end{tcolorbox}

\input{simulazioni/generati/simulazioni}
